% Stage 1a: the FPGA drives the Stage 0c breadboard; one lit LED walks along the 74HC595.
% Build with ./build.sh (figures, then latexmk; settings in .latexmkrc).

\documentclass[11pt,a4paper]{article}

\usepackage{fontspec}
\usepackage[margin=2.5cm]{geometry}
\usepackage{float}
\usepackage{enumitem}
\usepackage{booktabs}
\usepackage{minted}
\setminted{fontsize=\footnotesize, frame=single, framerule=0.4pt, rulecolor=black!30, framesep=4pt, xleftmargin=-8pt, xrightmargin=-8pt}
\usepackage{chapter}
\usepackage[colorlinks=true,linkcolor=blue!50!black,urlcolor=blue!50!black]{hyperref}

\setlength{\parindent}{0pt}
\setlength{\parskip}{0.6em}

\title{A walking bit on the \name{74HC595}}
\author{}
\date{}

\begin{document}
\maketitle

The Stage~0c breadboard, with the FPGA in place of the three buttons. It repeats what you did by hand: put one 1 into the shift register, then shift and latch once per tick, so a single lit LED walks from \name{QA} to \name{QH}. The FPGA feeds in a new 1 every eighth tick, so the walk repeats forever.

\section*{Rewiring}

\begin{enumerate}
  \item \textbf{Buttons out.} With the FPGA driving a line low, a pressed button would short its output straight to 3.3\,V. Take out the three buttons and the wires from them to the 3.3\,V rail.
  \item \textbf{The \name{SRCLK} capacitor out.} The FPGA’s edges do not bounce, and 100\,nF would slow them. The other 100\,nF, across the chip’s supply pins 16 and 8, stays.
  \item \textbf{The pull-downs stay.} The FPGA drives the lines much harder than 10\,kΩ pulls them, so they make no difference once it runs, and they hold the inputs low while it loads.
  \item \textbf{Power from the Tang.} The power module out; the board’s 3V3 pin to the $+$ rail, a GND pin to the $-$ rail. One supply and one ground for both. Eight LEDs at about 4\,mA each is 32\,mA, well within what the board’s regulator can supply.\here[check regulator]
  \item \textbf{Three wires.} Pin 25 to \name{SER} (14), pin 26 to \name{SRCLK} (11), pin 27 to \name{RCLK} (12). The \name{595} runs at 3.3\,V, so it needs 3.3\,V pins, and not every FPGA pin is: the LEDs’ and buttons’ pins run at 1.8\,V, and a 1.8\,V high is too low for the \name{595} to read as \name{HIGH}. The board’s schematic shows which are which: pins named \name{\_1V8} are 1.8\,V. Pins 25--30 are 3.3\,V and wired to nothing else on the board; 28--30 are left for the \name{74HC165} in Stage~2.
\end{enumerate}

\section*{Timing}

The \name{595} needs time between the steps you pressed by hand. Its datasheet gives the minimums at 2, 4.5 and 6\,V, but not at 3.3\,V; the 2\,V figures are the safe bound (\name{SN74HC595}, −40 to 85\,°C):

\begin{center}
\begin{tabular}{ll}
  \toprule
  Requirement & Minimum at 2\,V \\
  \midrule
  \name{SER} steady before \name{SRCLK} rises & 125\,ns \\
  \name{SRCLK} rises before \name{RCLK} rises & 94\,ns \\
  \name{SRCLK} or \name{RCLK} high, or low & 100\,ns \\
  \bottomrule
\end{tabular}
\end{center}

One cycle of the 27\,MHz clock is 37\,ns: too short. So the design spaces its steps 1\,µs apart (27 cycles), eight times the longest minimum. A whole tick, set \name{SER}, shift, latch, takes 3\,µs; the LEDs then hold still for the rest of the half second.

\section*{The design}

\name{rtl/walk.v}:

\begin{minted}{verilog}
module walk (
    input clk,              // 27 MHz
    output reg ser = 0,
    output reg srclk = 0,
    output reg rclk = 0
);

localparam TICK = 13_500_000;  // 0.5 s at 27 MHz
localparam STEP = 27;          // 1 µs

reg [23:0] t = 0;              // time into this tick
reg [2:0]  k = 0;              // which tick of eight: a 1 goes in on tick 0

always @(posedge clk) begin
    t <= (t == TICK - 1) ? 0 : t + 1;
    case (t)
        0:        ser <= (k == 0);
        STEP:     srclk <= 1;                   // shift
        2 * STEP: begin srclk <= 0; rclk <= 1; end  // latch
        3 * STEP: begin rclk <= 0; k <= k + 1; end
    endcase
end

endmodule
\end{minted}

The order is the one you pressed: \name{SER} first, then \name{SRCLK}, then \name{RCLK}. \texttt{k} counts to 7 and wraps to 0 on its own, being 3 bits wide, so every eighth tick sends in a new 1 just as the old one leaves \name{QH}.

\name{constraints/walk.cst}: \texttt{clk} on 52; \texttt{ser}, \texttt{srclk}, \texttt{rclk} on 25, 26, 27, all \texttt{LVCMOS33}. Then \texttt{make flash TOP=walk}, and \texttt{make schematic TOP=walk} for its RTL drawing: a counter, three comparators and four registers.\here[figure to come]

\section*{On the analyser}

After Stage~0b: CH0 on \name{SER}, CH1 on \name{SRCLK}, CH2 on \name{RCLK}, GND to GND. Trigger on CH1’s rising edge, at 4\,MS/s or more. Each tick is the same three edges, 1\,µs apart; \name{SER} is high on one tick in eight. CH3 on \name{QH'} (pin 9), the bit leaving the register, shows the 1 coming out eight ticks after it went in.

\section*{Done when}

One lit LED walks from \name{QA} to \name{QH}, a step every half second, and starts again at \name{QA}. After the first eight ticks, any pattern the chip powered up with has gone, and exactly one LED is lit at all times.

\end{document}
